% ============================================================
% COMANDI TIKZ PER COMPONENTI ELETTRONICI (versione corretta)
% ============================================================

% --- COLORI ---
\definecolor{bluscuro}{RGB}{20, 60, 130}
\definecolor{rossoscuro}{RGB}{170, 30, 30}
\definecolor{verdescuro}{RGB}{20, 120, 60}
\definecolor{arancione}{RGB}{230, 130, 20}
\definecolor{viola}{RGB}{120, 40, 160}

% --- Helper per l'etichetta ---
% #1 = angolo di rotazione del componente
% #2 = distanza dall'asse
% #3 = testo
\newcommand{\etichettaComponente}[3]{%
  \ifnum#1=0
    \node[above] at (0,#2) {#3};
  \else\ifnum#1=90
    \node[right] at (#2,0) {#3};
  \else\ifnum#1=-90
    \node[left] at (-#2,0) {#3};
  \else\ifnum#1=180
    \node[below] at (0,-#2) {#3};
  \else
    \node[above] at (0,#2) {#3};
  \fi\fi\fi\fi
}

% ============================================================
% RESISTORI
% ============================================================

% --- Resistore ANSI (a zig-zag) ---
\newcommand{\resistoreANSI}[4][0]{%
  \begin{scope}[shift={(#2,#3)}]
    \begin{scope}[rotate=#1]
      \draw[thick] (-1,0) -- (-0.7,0)
        -- (-0.6,0.3) -- (-0.4,-0.3) -- (-0.2,0.3) -- (0,-0.3)
        -- (0.2,0.3) -- (0.4,-0.3) -- (0.6,0.3) -- (0.7,0)
        -- (1,0);
    \end{scope}
    \etichettaComponente{#1}{0.55}{#4}
  \end{scope}
}

% --- Resistore IEC (rettangolo) ---
\newcommand{\resistoreIEC}[4][0]{%
  \begin{scope}[shift={(#2,#3)}]
    \begin{scope}[rotate=#1]
      \draw[thick] (-1,0) -- (-0.7,0);
      \draw[thick] (-0.7,-0.35) rectangle (0.7,0.35);
      \draw[thick] (0.7,0) -- (1,0);
    \end{scope}
    \etichettaComponente{#1}{0.55}{#4}
  \end{scope}
}

% ============================================================
% GENERATORI DI TENSIONE
% ============================================================

% --- Generatore "spiegazione" (con + e - colorati) ---
\newcommand{\generatoreSpiegazione}[4][0]{%
  \begin{scope}[shift={(#2,#3)}]
    \begin{scope}[rotate=#1]
      \draw[thick] (0,0) circle (0.5);
      \draw[thick] (0,-0.5) -- (0,0.5);
      \draw[thick] (0,0.5) -- (0,1);
      \draw[thick] (0,-0.5) -- (0,-1);
      \node at (-0.25,0.3) {\Huge\textbf{\textcolor{rossoscuro}{\(+\)}}};
      \node at (-0.25,-0.3) {\Huge\textbf{\textcolor{bluscuro}{\(-\)}}};
    \end{scope}
    \etichettaComponente{#1}{0.9}{#4}
  \end{scope}
}

% --- Generatore "normale" (solo +, cerchio semplice) ---
\newcommand{\generatore}[4][0]{%
  \begin{scope}[shift={(#2,#3)}]
    \begin{scope}[rotate=#1]
      \draw[thick] (0,0) circle (0.5);
      \node at (-0.25,0.7) {\Large\textbf{\textcolor{black}{\(+\)}}};
      \draw[thick] (0,0.5) -- (0,1);
      \draw[thick] (0,-0.5) -- (0,-1);
      \draw[thick] (0,-0.5) -- (0,0.5);
    \end{scope}
    \etichettaComponente{#1}{0.9}{#4}
  \end{scope}
}

% ============================================================
% CONDENSATORE
% ============================================================
\newcommand{\condensatore}[4][0]{%
  \begin{scope}[shift={(#2,#3)}]
    \begin{scope}[rotate=#1]
      \draw[thick] (-0.15,-0.5) -- (-0.15,0.5);
      \draw[thick] (0.15,-0.5) -- (0.15,0.5);
      \draw[thick] (-1,0) -- (-0.15,0);
      \draw[thick] (0.15,0) -- (1,0);
    \end{scope}
    \etichettaComponente{#1}{0.75}{#4}
  \end{scope}
}

% ============================================================
% INDUTTORE
% ============================================================
\newcommand{\induttore}[4][0]{%
  \begin{scope}[shift={(#2,#3)}]
    \begin{scope}[rotate=#1]
      \draw[thick] (-1,0) -- (-0.8,0)
        arc[radius=0.2, start angle=180, end angle=0]
        arc[radius=0.2, start angle=180, end angle=0]
        arc[radius=0.2, start angle=180, end angle=0]
        arc[radius=0.2, start angle=180, end angle=0]
        -- (1,0);
    \end{scope}
    \etichettaComponente{#1}{0.6}{#4}
  \end{scope}
}


% --- Massa (ground) ---
% Uso: \massa[rotazione]{x}{y}
%   - rotazione: opzionale, in gradi (default 0)
%   - x, y: coordinate del punto di connessione (dove arriva il filo)
% Il simbolo si sviluppa verso il basso (rotazione 0 = triangolo sotto)
\newcommand{\massa}[3][0]{%
  \begin{scope}[shift={(#2,#3)}, rotate=#1]
    % Terminale dal punto di connessione alla base del triangolo
    \draw[thick] (0,0) -- (0,-0.25);
    % Triangolo
    \draw[thick] (-0.3,-0.25) -- (0.3,-0.25) -- (0,-0.55) -- cycle;
  \end{scope}
}